% =====================================================================
%  Chapitre 1 -- Qu'est-ce qu'apprendre ? Demarche et generalisation
%  Partie a : definitions, familles, pipeline, surapprentissage
% =====================================================================
\section{Chapitre 1 --- Démarche et généralisation}

\chapterslide{Chapitre 1}{Qu'est-ce qu'apprendre ?}%
{6 h --- 4 h CM/TD + 2 h TP}%
{\textbf{Objectifs.} Définir l'apprentissage automatique ; distinguer supervisé,
non supervisé et renforcement ; décrire le pipeline complet ; reconnaître le
sous-apprentissage et le surapprentissage sur une courbe.}

\begin{frame}{Le problème, en une phrase}
\begin{columns}[t]
\begin{column}{0.48\textwidth}
\begin{wbox}[Ce qu'on ne veut plus faire]
Écrire les règles à la main. Pour trier un courriel en spam, il faudrait
lister tous les mots suspects, toutes les fautes, toutes les astuces.
Chaque nouvelle ruse casse la règle précédente.
\end{wbox}
\end{column}
\begin{column}{0.48\textwidth}
\begin{kbox}[Ce qu'on veut]
Donner des \textbf{exemples} à un programme et le laisser trouver la règle.
Il s'améliore avec l'expérience, sur une tâche, mesurée par une
\textbf{performance}.
\end{kbox}
\end{column}
\end{columns}
\vspace{2mm}
\begin{solbox}[Définition (Mitchell, 1997)]
Un programme \emph{apprend} à partir de l'expérience $E$ pour une tâche $T$ et
une mesure de performance $P$ si sa performance sur $T$, mesurée par $P$,
s'améliore avec $E$.
\end{solbox}
\end{frame}

\begin{frame}{Apprendre, ce n'est pas mémoriser}
\begin{columns}[t]
\begin{column}{0.52\textwidth}
\begin{itemize}
  \item Mémoriser les 1\,000 exemples vus ne sert à rien le jour de l'examen.
  \item Ce qu'on veut : être bon sur des exemples \textbf{jamais vus}.
  \item C'est la \textbf{généralisation} --- et c'est toute la difficulté.
\end{itemize}
\begin{kbox}[La conséquence immédiate]
On ne peut pas juger un modèle sur les données qui ont servi à l'entraîner.
Il faut un jeu mis de côté.
\end{kbox}
\end{column}
\begin{column}{0.44\textwidth}
\begin{ebox}[L'analogie de l'étudiant]
Un étudiant qui apprend par cœur les annales et réussit le même sujet
n'a rien prouvé. On le juge sur un sujet nouveau.
\end{ebox}
\end{column}
\end{columns}
\end{frame}

\begin{frame}{Trois familles d'apprentissage}
\begin{center}
\includegraphics[width=0.97\textwidth]{g02_types.png}
\end{center}
\begin{columns}[t]
\begin{column}{0.32\textwidth}
\footnotesize\textbf{Supervisé} : on connaît la réponse. Classification,
régression. C'est 80 \% de ce cours.
\end{column}
\begin{column}{0.32\textwidth}
\footnotesize\textbf{Non supervisé} : aucune étiquette. Regrouper, réduire,
détecter l'anomalie.
\end{column}
\begin{column}{0.32\textwidth}
\footnotesize\textbf{Renforcé} : apprendre par essai-erreur, avec une
récompense différée.
\end{column}
\end{columns}
\end{frame}

\begin{frame}{Des exemples qui parlent}
\begin{center}
\small
\begin{tabular}{@{}llll@{}}
\toprule
\textbf{Tâche} & \textbf{Type} & \textbf{Entrée} & \textbf{Sortie} \\
\midrule
Filtrer le spam & supervisé, classification & mots du courriel & spam / non spam \\
Prédire un loyer & supervisé, régression & surface, quartier & montant en FCFA \\
Segmenter des clients & non supervisé & achats, fréquence & groupes \\
Détecter une panne & supervisé (série) & métriques serveur & alerte \\
Jouer à un jeu & renforcé & état du jeu & action \\
\bottomrule
\end{tabular}
\end{center}
\vspace{1mm}
\begin{kbox}[Le réflexe à prendre]
Avant tout modèle : \textbf{quelle est la sortie attendue ?} Un nombre
(régression), une classe (classification), un groupe (non supervisé) ?
Tout le reste en découle.
\end{kbox}
\end{frame}

\begin{frame}{Le pipeline, et l'endroit où tout casse}
\begin{center}
\includegraphics[width=0.98\textwidth]{g01_pipeline.png}
\end{center}
\vspace{1mm}
\begin{wbox}[Le modèle n'est pas le cœur du problème]
Dans la pratique, les erreurs viennent surtout des étapes 1, 2 et 4 :
données mal comprises, variables construites avec fuite, évaluation trompeuse.
Un modèle simple bien évalué bat un modèle complexe mal évalué.
\end{wbox}
\end{frame}

\begin{frame}{Le vocabulaire minimal}
\begin{center}
\small
\begin{tabular}{@{}lll@{}}
\toprule
\textbf{Notation} & \textbf{Nom} & \textbf{Exemple} \\
\midrule
$X$ & matrice des variables & 569 patients $\times$ 30 mesures \\
$y$ & étiquette / cible & tumeur bénigne ou maligne \\
$f_\theta$ & modèle de paramètres $\theta$ & régression logistique \\
$L(y, f_\theta(x))$ & fonction de perte & entropie croisée \\
$P$ & métrique de performance & exactitude, rappel, RMSE \\
\bottomrule
\end{tabular}
\end{center}
\vspace{1.5mm}
\begin{kbox}[Paramètres vs hyperparamètres]
\textbf{Paramètres} : appris par l'optimisation (poids $\theta$).
\textbf{Hyperparamètres} : choisis par vous (degré du polynôme, profondeur
d'un arbre, $k$ du k-means). Ils se règlent sur la \textbf{validation}, jamais
sur le test.
\end{kbox}
\end{frame}

\begin{frame}{Perte et métrique : ne pas confondre}
\begin{columns}[t]
\begin{column}{0.48\textwidth}
\begin{ebox}[La perte --- ce que le modèle optimise]
Différentiable, calculée sur chaque exemple : erreur quadratique,
entropie croisée. C'est le moteur de l'entraînement.
\end{ebox}
\end{column}
\begin{column}{0.48\textwidth}
\begin{ebox}[La métrique --- ce que vous décidez]
Ce qui compte pour l'utilisateur : RMSE en FCFA, rappel sur les pannes,
F1 sur les cas rares. Elle n'est pas toujours différentiable.
\end{ebox}
\end{column}
\end{columns}
\vspace{2mm}
\begin{wbox}[L'erreur classique]
Minimiser une perte qui ne correspond pas à la métrique métier : un modèle
qui optimise l'exactitude sur des données déséquilibrées prédit toujours
la classe majoritaire.
\end{wbox}
\end{frame}

\begin{frame}{Le surapprentissage, visible sur une courbe}
\begin{center}
\includegraphics[width=0.56\textwidth]{g03_surapprentissage.png}
\end{center}
\vspace{1mm}
\begin{solbox}[Lecture (calculée, graine 4112)]
Degré 5 : entraînement $0{,}042$, test $0{,}052$. Plus loin, l'entraînement
descend encore mais le test remonte : le modèle apprend le \textbf{bruit}.
\end{solbox}
\end{frame}

\begin{frame}{Trois régimes, trois remèdes}
\begin{center}
\small
\footnotesize
\begin{tabular}{@{}p{0.24\textwidth}p{0.28\textwidth}p{0.38\textwidth}@{}}
\toprule
\textbf{Régime} & \textbf{Symptôme} & \textbf{Remède} \\
\midrule
Sous-apprentissage & entraînement \textbf{élevé} & modèle plus complexe,
meilleures variables \\
Bon compromis & les deux erreurs basses & --- \\
Surapprentissage & entraînement \textbf{très bas}, test \textbf{élevé} &
plus de données, régularisation, modèle plus simple \\
\bottomrule
\end{tabular}
\normalsize
\end{center}
\vspace{2mm}
\begin{kbox}[La question à poser devant toute courbe]
L'écart entraînement/validation se creuse-t-il ? Si oui, le problème n'est pas
dans les données mais dans la \textbf{capacité du modèle}.
\end{kbox}
\end{frame}

\begin{frame}{Quiz 1 --- fondements}
\begin{columns}[t]
\begin{column}{0.49\textwidth}
\begin{qbox}[Question 1]
Un modèle obtient 99 \% en entraînement et 61 \% en test. Que se passe-t-il ?
\begin{itemize}
  \item[(a)] sous-apprentissage \quad (b) surapprentissage
  \item[(c)] fuite de données
\end{itemize}
\end{qbox}
\begin{qbox}[Question 2]
La différence entre paramètre et hyperparamètre :
\begin{itemize}
  \item[(a)] le paramètre est appris, l'hyperparamètre est choisi
  \item[(b)] l'hyperparamètre est plus important
  \item[(c)] aucune, ce sont des synonymes
\end{itemize}
\end{qbox}
\end{column}
\begin{column}{0.49\textwidth}
\begin{qbox}[Question 3]
Pourquoi ne pas juger un modèle sur ses données d'entraînement ?
\begin{itemize}
  \item[(a)] c'est plus lent
  \item[(b)] cela ne mesure pas la généralisation
  \item[(c)] les données sont trop nombreuses
\end{itemize}
\end{qbox}
\begin{solbox}[Réponses]
\textbf{1.} (b). \quad \textbf{2.} (a). \quad \textbf{3.} (b) : on mesure la
mémorisation, pas la capacité à traiter un cas nouveau.
\end{solbox}
\end{column}
\end{columns}
\end{frame}
